\documentclass[11pt]{article}
\usepackage[margin=0.9in]{geometry}
\usepackage{amsmath,amssymb,amsthm,booktabs,mathrsfs}
\usepackage[T1]{fontenc}
\usepackage{lmodern,microtype}
\usepackage[colorlinks=true,linkcolor=blue,urlcolor=blue]{hyperref}
\newtheorem{theorem}{Theorem}[section]
\newtheorem{lemma}[theorem]{Lemma}
\newtheorem{proposition}[theorem]{Proposition}
\newtheorem{corollary}[theorem]{Corollary}
\theoremstyle{remark}
\newtheorem{remark}[theorem]{Remark}
\newcommand{\norm}[1]{\left\|#1\right\|}
\title{Toward a four-dimensional Yang--Mills construction\\
\large Notebook 22: transport into the next coarse background}
\author{Research working document}
\date{17 October 2026}
\begin{document}
\maketitle
\begin{abstract}
We construct a strictly supported expansion of the actual physical
stationary current background as a function of the next-coarse
field, uniformly in preceding classical refinement and periodic
volume. An affine version of the exact ordinary pivot chart and
the retained flat inverse yield a support-valued stationary
equation. Coordinatewise Cauchy estimates convert the previously
localized scalar action into an analytic gradient majorant without
a factorial loss. The solution is identified with the established
gauge-zero classical branch. We then transport holomorphic
connected-carrier activities through that branch with explicit
losses of analytic radius and support weight. The number of
Taylor coordinates is counted exactly, and the vacuum term is
separated. This supplies a bounded one-step transport for the
generated quantum activities. It does not assert contraction
or stability of the analytic domains under repeated integration.
\end{abstract}

Notebook 21 localized the classical action and controlled one
integration of it in independent current backgrounds. A new
coarse effective action requires the substitution of its
stationary current background. The following two results prove
that substitution and its activity norm estimate. Their domain
and weight losses are retained as inputs to the remaining
quantum induction.

